% GENERATED FILE. Edit canonical lessons, then run npm run generate:books.
% canonical-source-hash: fnv1a64:ed957eeaf76e1552
% canonical-lessons: TE-C354-kutra, TE-C354-mabhyapettu, TE-C354-bayatapettu, TE-C354-kappipuccu, TE-C354-gudhacari

\chapter{A Plot Uncovered}
\label{ch:te-a-plot-uncovered}
\hlchaptermodality{354}

\begin{chapteropening}
\textbf{By the end of this chapter:} \emph{I can say a conspiracy, to mislead, to expose, to cover up and a spy.}

Complete the last lesson of chapter 354: I can say a conspiracy, to mislead, to expose, to cover up and a spy.
\end{chapteropening}

\section[\texorpdfstring{\te{కుట్ర} (kuṭra)}{కుట్ర (kuṭra)}]{\te{కుట్ర} (kuṭra) — a conspiracy; a plot}

\label{lesson:TE-C354-kutra}



\begin{quote}
Before the new one: say the Telugu for to appear suddenly, then the Telugu for to detect.
\end{quote}

\subsection*{You'll want to know: \te{కుట్ర}}
\textbf{\te{కుట్ర}} — \emph{kuṭra} — \textquotedblleft{}a conspiracy; a plot\textquotedblright{}.

\begin{quote}
\te{రాజు} \te{మీద} \te{ఒక} \te{కుట్ర} \te{జరిగింది}.
\end{quote}

\begin{quote}
\emph{rāju mīda oka kuṭra jarigindi.}
\end{quote}

\begin{quote}
There was a plot against the king.
\end{quote}

\te{కుట్ర} \te{పన్ను}, to hatch a plot.

\begin{cousinweb}
\textbf{\te{కుట్ర}} (\emph{kuṭra}) is a Telugu word for a secret plot; its origin is not settled.
\end{cousinweb}

\begin{grammarlens}[title={What you've built}]
One more word for telling what happened.
\end{grammarlens}

\subsection*{Your turn}
\begin{itemize}
\raggedright
  \item \emph{Say it:} \emph{kuṭra}
  \item \emph{Say it:} \emph{kuṭra}, once more
  \item \emph{Say it:} \emph{rāju mīda oka kuṭra jarigindi.} — There was a plot against the king.
  \item \emph{Recall:} say the Telugu for to appear suddenly, then the Telugu for to detect, then say \emph{kuṭra} again
\end{itemize}

\begin{tcolorbox}[breakable,colback=teal!4,colframe=teal!35!black,title={Before you move on}]
What does \te{కుట్ర} mean? (\textquotedblleft{}A conspiracy; a plot\textquotedblright{}.) Say it once more.
\end{tcolorbox}
\section[\texorpdfstring{\te{మభ్యపెట్టు} (mabhyapeṭṭu)}{మభ్యపెట్టు (mabhyapeṭṭu)}]{\te{మభ్యపెట్టు} (mabhyapeṭṭu) — to mislead; to hoodwink}

\label{lesson:TE-C354-mabhyapettu}



\begin{quote}
Before the new one: say the Telugu for to detect, then the Telugu for a conspiracy.
\end{quote}

\subsection*{You'll want to know: \te{మభ్యపెట్టు}}
\textbf{\te{మభ్యపెట్టు}} — \emph{mabhyapeṭṭu} — \textquotedblleft{}to mislead; to hoodwink\textquotedblright{}.

\begin{quote}
\te{అతను} \te{అబద్ధాలతో} \te{అందరినీ} \te{మభ్యపెట్టాడు}.
\end{quote}

\begin{quote}
\emph{atanu abaddhālatō andarinī mabhyapeṭṭāḍu.}
\end{quote}

\begin{quote}
He misled everyone with lies.
\end{quote}

Past (I …): \te{మభ్యపెట్టాను}, \emph{mabhyapeṭṭānu}.

To cloud someone's mind; a written-style word.

\begin{cousinweb}
\textbf{\te{మభ్యపెట్టు}} (\emph{mabhyapeṭṭu}) ends in native \emph{peṭṭu}, \textquotedblleft{}to put\textquotedblright{} (DEDR 4390); the origin of the first part is not settled.
\end{cousinweb}

\begin{grammarlens}[title={What you've built}]
One more word for telling what happened.
\end{grammarlens}

\subsection*{Your turn}
\begin{itemize}
\raggedright
  \item \emph{Say it:} \emph{mabhyapeṭṭu}
  \item \emph{Say it:} \emph{mabhyapeṭṭu}, once more
  \item \emph{Say it:} \emph{atanu abaddhālatō andarinī mabhyapeṭṭāḍu.} — He misled everyone with lies.
  \item \emph{Recall:} say the Telugu for to detect, then the Telugu for a conspiracy, then say \emph{mabhyapeṭṭu} again
\end{itemize}

\begin{tcolorbox}[breakable,colback=teal!4,colframe=teal!35!black,title={Before you move on}]
What does \te{మభ్యపెట్టు} mean? (\textquotedblleft{}To mislead; to hoodwink\textquotedblright{}.) Say it once more.
\end{tcolorbox}
\section[\texorpdfstring{\te{బయటపెట్టు} (bayaṭapeṭṭu)}{బయటపెట్టు (bayaṭapeṭṭu)}]{\te{బయటపెట్టు} (bayaṭapeṭṭu) — to expose; to bring out}

\label{lesson:TE-C354-bayatapettu}



\begin{quote}
Before the new one: say the Telugu for a conspiracy, then the Telugu for to mislead.
\end{quote}

\subsection*{You'll want to know: \te{బయటపెట్టు}}
\textbf{\te{బయటపెట్టు}} — \emph{bayaṭapeṭṭu} — \textquotedblleft{}to expose; to bring out\textquotedblright{}.

\begin{quote}
\te{విలేకరి} \te{ఆ} \te{కుట్రను} \te{బయటపెట్టాడు}.
\end{quote}

\begin{quote}
\emph{vilēkari ā kuṭranu bayaṭapeṭṭāḍu.}
\end{quote}

\begin{quote}
The reporter exposed the plot.
\end{quote}

Past (I …): \te{బయటపెట్టాను}, \emph{bayaṭapeṭṭānu}.

\te{బయట}, outside, with \te{పెట్టు}, to put: to put a secret outside.

\begin{cousinweb}
\textbf{\te{బయటపెట్టు}} (\emph{bayaṭapeṭṭu}) is a native compound: \emph{bayaṭa}, \textquotedblleft{}outside\textquotedblright{}, which sits beside \emph{bayalu}, \textquotedblleft{}the outside\textquotedblright{} (DEDR 3940), and \emph{peṭṭu}, \textquotedblleft{}to put\textquotedblright{} (DEDR 4390).
\end{cousinweb}

\begin{grammarlens}[title={What you've built}]
One more word for telling what happened.
\end{grammarlens}

\subsection*{Your turn}
\begin{itemize}
\raggedright
  \item \emph{Say it:} \emph{bayaṭapeṭṭu}
  \item \emph{Say it:} \emph{bayaṭapeṭṭu}, once more
  \item \emph{Say it:} \emph{vilēkari ā kuṭranu bayaṭapeṭṭāḍu.} — The reporter exposed the plot.
  \item \emph{Recall:} say the Telugu for a conspiracy, then the Telugu for to mislead, then say \emph{bayaṭapeṭṭu} again
\end{itemize}

\begin{tcolorbox}[breakable,colback=teal!4,colframe=teal!35!black,title={Before you move on}]
What does \te{బయటపెట్టు} mean? (\textquotedblleft{}To expose; to bring out\textquotedblright{}.) Say it once more.
\end{tcolorbox}
\section[\texorpdfstring{\te{కప్పిపుచ్చు} (kappipuccu)}{కప్పిపుచ్చు (kappipuccu)}]{\te{కప్పిపుచ్చు} (kappipuccu) — to cover up; to hush up}

\label{lesson:TE-C354-kappipuccu}



\begin{quote}
Before the new one: say the Telugu for to mislead, then the Telugu for to expose.
\end{quote}

\subsection*{You'll want to know: \te{కప్పిపుచ్చు}}
\textbf{\te{కప్పిపుచ్చు}} — \emph{kappipuccu} — \textquotedblleft{}to cover up; to hush up\textquotedblright{}.

\begin{quote}
\te{అధికారులు} \te{ఆ} \te{తప్పును} \te{కప్పిపుచ్చారు}.
\end{quote}

\begin{quote}
\emph{adhikārulu ā tappunu kappipuccāru.}
\end{quote}

\begin{quote}
The officials covered up the mistake.
\end{quote}

Past (I …): \te{కప్పిపుచ్చాను}, \emph{kappipuccānu}.

Built on \te{కప్పు}, to cover.

\begin{cousinweb}
\textbf{\te{కప్పిపుచ్చు}} (\emph{kappipuccu}) is a native compound built on \emph{kappu}, \textquotedblleft{}to cover, conceal\textquotedblright{} (DEDR 1221), which also gives \te{కప్పు}, a roof.
\end{cousinweb}

\begin{grammarlens}[title={What you've built}]
One more word for telling what happened.
\end{grammarlens}

\subsection*{Your turn}
\begin{itemize}
\raggedright
  \item \emph{Say it:} \emph{kappipuccu}
  \item \emph{Say it:} \emph{kappipuccu}, once more
  \item \emph{Say it:} \emph{adhikārulu ā tappunu kappipuccāru.} — The officials covered up the mistake.
  \item \emph{Recall:} say the Telugu for to mislead, then the Telugu for to expose, then say \emph{kappipuccu} again
\end{itemize}

\begin{tcolorbox}[breakable,colback=teal!4,colframe=teal!35!black,title={Before you move on}]
What does \te{కప్పిపుచ్చు} mean? (\textquotedblleft{}To cover up; to hush up\textquotedblright{}.) Say it once more.
\end{tcolorbox}
\section[\texorpdfstring{\te{గూఢచారి} (gūḍhacāri)}{గూఢచారి (gūḍhacāri)}]{\te{గూఢచారి} (gūḍhacāri) — a spy}

\label{lesson:TE-C354-gudhacari}



\begin{quote}
Before the new one: say the Telugu for to expose, then the Telugu for to cover up.
\end{quote}

\subsection*{You'll want to know: \te{గూఢచారి}}
\textbf{\te{గూఢచారి}} — \emph{gūḍhacāri} — \textquotedblleft{}a spy\textquotedblright{}.

\begin{quote}
\te{శత్రువు} \te{గూఢచారి} \te{కోటలోకి} \te{వచ్చాడు}.
\end{quote}

\begin{quote}
\emph{śatruvu gūḍhacāri kōṭalōki vaccāḍu.}
\end{quote}

\begin{quote}
An enemy spy came into the fort.
\end{quote}

Plural: \te{గూఢచారులు}, \emph{gūḍhacārulu}.

\te{గూఢ}, hidden, and \te{చారి}, one who goes about.

\begin{cousinweb}
\textbf{\te{గూఢచారి}} (\emph{gūḍhacāri}) is borrowed from Sanskrit: \emph{gūḍha}, \textquotedblleft{}hidden\textquotedblright{} (CDIAL 4223), and \emph{cārin}, \textquotedblleft{}moving; a spy\textquotedblright{} (CDIAL 4765).
\end{cousinweb}

\begin{grammarlens}[title={What you've built}]
One more word for telling what happened.
\end{grammarlens}

\subsection*{Your turn}
\begin{itemize}
\raggedright
  \item \emph{Say it:} \emph{gūḍhacāri}
  \item \emph{Say it:} \emph{gūḍhacāri}, once more
  \item \emph{Say it:} \emph{śatruvu gūḍhacāri kōṭalōki vaccāḍu.} — An enemy spy came into the fort.
  \item \emph{Recall:} say the Telugu for to expose, then the Telugu for to cover up, then say \emph{gūḍhacāri} again
\end{itemize}

\begin{tcolorbox}[breakable,colback=teal!4,colframe=teal!35!black,title={Before you move on}]
What does \te{గూఢచారి} mean? (\textquotedblleft{}A spy\textquotedblright{}.) Say it once more.
\end{tcolorbox}
